\subsection{The original character data satisfy the construction hypotheses}
\label{sec:support-auxiliary-admission}

We check the two families of characters used in the arithmetic
argument.  In both, the character entering the deformation problem
is the compensated character $\Psi=\varepsilon^e\chi$, where
$\chi$ is the given finite Dirichlet character.  Its parity, its
prime-power components and its residual character must all be kept.

\begin{lemma}[A character of large $2$-power conductor]
\label{lem:support-dyadic-character}
Let $\alpha$ be a character of $\mathbb Z_2^\times$ of primitive
conductor $2^c$, with $c\geq4$.  If $p>1$ is odd and $o\geq1$
satisfies $8op^2<2^c$, then $\alpha(p)^{4o}\ne1$.
There is also $u\equiv-1\pmod4$ with $\alpha(u)=1$.
\end{lemma}

\begin{proof}
On $1+4\mathbb Z_2$, the character has kernel exactly
$1+2^c\mathbb Z_2$.  Indeed $1+4\mathbb Z_2$ is procyclic
with generator $5$, and a smaller order for $\alpha(5)$
would lower its conductor.  Write
$s=v_2(p^2-1)$ and $r=v_2(o)$.  Repeatedly factoring
$x^2-1$ gives
\[
 v_2(p^{4o}-1)=s+r+1,
 \qquad 2^{s+r+1}\leq2o(p^2-1)<8op^2<2^c.
\]
Thus $p^{4o}$ is outside the kernel on $1+4\mathbb Z_2$.
For the last assertion, $\alpha(-1)\in\{\pm1\}$ lies in
the image of $1+4\mathbb Z_2$, whose image is cyclic of
order $2^{c-2}\geq4$.  Choose $v\in1+4\mathbb Z_2$ with
$\alpha(v)=\alpha(-1)^{-1}$ and take $u=-v$.
\end{proof}

\begin{proposition}[The family with quadratic field $\mathbb Q(i)$]
\label{prop:support-host-admission}
Let $\ell$ be odd and let $\chi$ be the original finite character
of modulus $N$, with $\Psi=\varepsilon^e\chi$ even.  Write
$\chi^{(0)}$ and $\chi^{(\ell)}$ for its parts of order
prime to $\ell$ and of $\ell$-power order, respectively.
Let $2^c$ be the conductor of its component at $2$.
Assume $c\geq4$ and, for the given $o\geq1$,
\[
  \operatorname{ord}(\chi^{(0)}_p)\mid2o\quad
       \text{for every odd prime }p,
  \qquad 8o\ell^2<2^c.                                   \tag{AA1}
\]
Then the hypotheses of Theorem~\ref{thm:support-auxiliary-tower}
hold for the original $\eta=\eta_4$, $M=\ell N$, and the
field $F_b$ cut out by the product of the components of
$\chi^{(\ell)}$ away from $\ell$.
\end{proposition}

\begin{proof}
The field $F_b$ is abelian of $\ell$-power degree and
unramified at $\ell$ by its definition.  At an inertia
group above $p\ne\ell$, the cyclotomic character and the
$\ell$-component of the finite character are trivial.
On $G_{F_b}$ the remaining $\ell$-power part is trivial
as well.  Hence $\Psi$ on this inertia is a root of unity
of order prime to $\ell$.  Such a root whose residue is
$1$ equals $1$; reduction is injective on roots of unity
of order prime to $\ell$.  This proves the precise
condition on $\Psi$ away from $\ell$.

The characters $\eta_4$ and $\overline\Psi$ factor through
roots of unity of order dividing a power of $M$.
The field $\mathbb Q(i)$ is unramified at odd $\ell$,
and $\overline\Psi$ is even.  Its inertia character at
$\ell$ is a power of $\omega$: wild inertia has pro-$\ell$
image and hence trivial image in $k^\times$, while the
tame abelian quotient of inertia maps through
$\mathbb F_\ell^\times$.

Use the last assertion of Lemma~\ref{lem:support-dyadic-character}
for the component $\chi_2$, whose order is prime to $\ell$.
Local cyclotomic surjectivity at $2$ gives inertia with
this value $u$ on the $2$-power roots of unity.  Its
values on all odd-primary character components and on
$\omega$ are $1$.  Hence $\overline\Psi=\omega=1$ there
while $\eta_4=-1$.  This proves the required failure of
the kernel inclusion.

Finally choose a Frobenius lift $s\in D_B$ at $\ell$
which has $\omega(s)=1$.  Such a lift exists because
$\mathbb Q_\ell(\mu_\ell)$ is totally ramified, and its
intersection with the maximal unramified extension is
trivial.  On roots of unity of order prime to $\ell$,
$s$ acts by the $\ell$th power.  Its value on the
prime-to-$\ell$ part of the local character at $\ell$
is $1$, since that character factors through $\omega$.
By (AA1), raising the other odd-prime components of
$\chi^{(0)}(s)$ to $4o$ kills them all.  Consequently
\[
     \chi^{(0)}(s)^{4o}=\chi_2(\ell)^{4o}\ne1.
\]
Reduction preserves this inequality, and kills the
$\ell$-power part of $\chi$.  Thus
$\overline\Psi(s)^4\ne1$.  For $t=s^2$ we have
$\eta_4(t)=1$ and $\overline\Psi(t)^2\ne1$, exactly the
local condition needed in the tower construction.
\end{proof}

\begin{proposition}[The original quadratic character data]
\label{prop:support-good-admission}
Let $K$ be the original imaginary quadratic field, unramified
at $\ell$, and let $\eta$ be its character.  In the second
family, the original finite character $\chi$ has order
prime to $\ell$, its local component at $\ell$ has square
$1$, and the compensated character $\Psi=\varepsilon^e\chi$
is even.  Retain the original conditions
\[
 \ell-1\nmid2e,
 \qquad \ker\omega\cap\ker\chi\nsubseteq\ker\eta.      \tag{AA2}
\]
Then the hypotheses of Theorem~\ref{thm:support-auxiliary-tower}
hold with the original $K$, $M=\ell N|\operatorname{disc}(K)|$
and $F_b=\mathbb Q$.
\end{proposition}

\begin{proof}
Away from $\ell$, $\Psi$ is a root of unity of order
prime to $\ell$, so the required implication from
residual triviality to triviality follows by injectivity
of reduction.  The cyclotomic and conductor assertions
are immediate from the given moduli.  In particular the
negative quadratic discriminant supplies $D>0$ with
$K=\mathbb Q(\sqrt{-D})$ and all its prime factors in $M$.
The original modulus is even, so $2\ell\mid M$.

The local character of $\chi$ at $\ell$ is a power
$\omega^a$ after reduction.  Its square being trivial
means $\ell-1\mid2a$.  Thus
$\overline\Psi|_{I_\ell}=\omega^{e+a}$ has nontrivial
square by (AA2).  Choose inertia where its square is
nontrivial; $\eta=1$ on that inertia because $K$ is
unramified at $\ell$.  This gives the required local
element.  Reduction is injective on the values of
$\chi$, so on $\ker\omega$ the kernels of $\chi$ and
$\overline\Psi$ agree.  The second condition in (AA2)
is therefore precisely the required kernel condition.
The parity condition is the parity of $\Psi$ itself.
All hypotheses have now been checked on the original
fields and characters.
\end{proof}

In the first family, the original conductor condition applies
at every prime in its prescribed set, which includes $\ell$;
only that instance was used in (AA1).  In the second family,
the character away from $\ell$ used in the original kernel
condition differs from $\chi$ by a power of $\omega$, so
their kernels agree after intersection with $\ker\omega$.
Consequently neither proposition strengthens the original
input hypotheses\cite{AuxTowerSource}.  Together with
Theorems~\ref{thm:support-auxiliary-tower} and
\ref{thm:support-auxiliary-frame} and
Proposition~\ref{prop:support-auxiliary-seed}, they supply
the complete original auxiliary field, Galois quotient,
cohomology, degree and seed conditions in both families.
The fixed-index subgroup in Theorem~\ref{thm:support-finite-end}
can therefore be chosen from these constructed arithmetic
data; its degree is still allowed to grow with the original
parameters.
